\section{Evaluation --- draft v1}
\label{sec:eval}



\subsection{Methodology}

We evaluate three measurable questions from \S3: the invasiveness of
the mechanisms (H1/E5), their always-on runtime and memory cost
(H3/E1/E2), and the structural-quality properties of the diagnostics
they enable (E3/E4). Human benefit is explicitly \emph{not} evaluated
(\S13 designs that study; \S11 explains why the claim is withheld).

\textbf{Workload.} Deterministic, seeded, \emph{valid} programs at 100 / 1,000
/ 5,000 / 20,000 lines (many 20--40-statement functions of mixed-depth
integer arithmetic). Valid programs are the fair worst case for this
architecture: every mechanism pays its cost and no diagnostic ever
redeems it.

\textbf{Measurement path.} \texttt{--check}: every stage through assembly text,
no JSON serialization, no toolchain subprocess, nothing printed on
success --- the stopwatch sees the compiler and only the compiler. This
mode exists because our first campaign timed \texttt{--json} and measured
the harness instead (\S9.2, "a discarded campaign").

\textbf{Design.} Factorial ablation over two compile-time flags
(recording, provenance), giving four binaries: FULL, NOTRACE,
NOPROV, BASELINE --- each mechanism's cost is attributed separately.
Ablation is by macro so argument expressions vanish too; a no-op
class would understate the cost it exists to measure (\S6.3, D13).

\textbf{Protocol.} Per (config, size): 3 warmups, 30 timed runs, 5
peak-working-set runs (sampled live; the counter is unreadable after
exit). Binaries built \texttt{-O2}. Statistics: median and MAD per cell ---
wall times on a desktop OS are right-skewed by scheduler and
antivirus interference, and medians resist those spikes; overhead is
(median\_config -- median\_baseline)/median\_baseline with a 95\%
bootstrap CI (10,000 resamples, percentile method, fixed seed).
Alternatives rejected: mean$\pm$SD assumes symmetry the data does not
have; comparing against production-compiler \emph{speed} would be
meaningless at this scale, so all comparisons are self-vs-self.

\textbf{Environment.} Intel Core i7-8665U @ 1.90 GHz, 32 GiB RAM,
Windows 11 Pro (build 26200), GCC 6.3.0 (MinGW.org, 32-bit target).
Single machine, single toolchain: an external-validity limit stated
in \S11, not an inconvenience to hide.

\subsection{Overhead (E1/E2)}

\textbf{A discarded campaign (methodological note).} Our first E1
collection timed \texttt{--json} with output piped through the measurement
shell. The numbers failed the coherence check we ran before accepting
any of them: ablated configurations measured \emph{slower} than the full
build, orderings flipped between sizes, and "baseline" was nine
seconds for a 20k-line file through an O(n) pipeline. The harness was
ingesting megabytes of serialized pipeline state per run --- we had
measured PowerShell. The dataset was deleted, the \texttt{--check} path
built, and the campaign rerun. We report this because measurement
pitfalls that are survived silently get repeated by readers.

**Naive recorder (eager detail strings) --- the result that refuted
H3.** Median $\pm$ MAD, ms:

\begin{table}[t]
\centering
\small
\begin{tabular}{lllll}
\toprule
lines & baseline & notrace & noprov & full \\
\midrule
100 & 46.9 $\pm$ 2.1 & 53.4 $\pm$ 2.5 & 58.3 $\pm$ 3.4 & 73.8 $\pm$ 4.1 \\
1,000 & 154.8 $\pm$ 14.9 & 188.7 $\pm$ 8.8 & 278.4 $\pm$ 5.2 & 335.6 $\pm$ 11.4 \\
5,000 & 570.9 $\pm$ 7.1 & 824.3 $\pm$ 48.3 & 1,569.9 $\pm$ 59.0 & 1,493.6 $\pm$ 53.5 \\
20,000 & 2,266.0 $\pm$ 133.0 & 2,237.9 $\pm$ 67.5 & 5,498.2 $\pm$ 312.7 & 5,824.2 $\pm$ 221.1 \\
\bottomrule
\end{tabular}
\end{table}

At 20,000 lines --- the startup-amortized row --- trace recording cost
+143\% [130, 168] (NOPROV) and the full configuration +157\%
[144, 204]. The cause is identifiable, not mysterious: the lexer
opens two scopes per token, each eagerly building detail strings ---
order 10$^5$--10$^6$ heap constructions per clean compile, all discarded.
This campaign also \emph{appeared} to show provenance as free (NOTRACE vs
BASELINE: --1.2\%, CI spanning zero) --- a reading the interleaved
protocol below corrects; under blocked measurement, whichever config
runs nearest the baseline in time inherits flattering drift.

\textbf{A second protocol iteration.} Re-measuring the lazy recorder under
the same blocked protocol produced medians that *decrease
monotonically in run order* across configs (1890 $\rightarrow$ 1813 $\rightarrow$ 1532 $\rightarrow$
1512 ms at 20k) --- the machine sped up as the campaign progressed, and
with mechanism costs now small, that drift dominated the
between-config deltas. The final protocol therefore interleaves
configs round-robin within each run index, rotating cycle order to
cancel position bias. Blocked data is archived, not used.

**Lazy recorder (D14), interleaved protocol --- the authoritative
table.** Fig. F7 plots both campaigns' overhead on one shared scale
(the flat right panel is the repair); the figure is generated from
the raw CSVs by a script that hard-asserts every plotted median
against the tables below, so figure and paper cannot silently
diverge. Median $\pm$ MAD, ms:

\begin{table}[t]
\centering
\small
\begin{tabular}{lllll}
\toprule
lines & baseline & notrace & noprov & full \\
\midrule
100 & 24.1 $\pm$ 0.9 & 27.0 $\pm$ 1.9 & 24.4 $\pm$ 1.0 & 25.8 $\pm$ 0.8 \\
1,000 & 78.6 $\pm$ 2.5 & 89.9 $\pm$ 2.6 & 79.9 $\pm$ 2.8 & 90.3 $\pm$ 2.5 \\
5,000 & 398.4 $\pm$ 12.1 & 467.2 $\pm$ 11.0 & 402.7 $\pm$ 12.1 & 467.8 $\pm$ 10.2 \\
20,000 & 1,549.5 $\pm$ 76.1 & 1,801.6 $\pm$ 76.5 & 1,585.8 $\pm$ 94.7 & 1,844.5 $\pm$ 94.2 \\
\bottomrule
\end{tabular}
\end{table}

Overhead vs baseline (95\% bootstrap CI):

\begin{table}[t]
\centering
\small
\begin{tabular}{llll}
\toprule
lines & full (total) & trace only (NOPROV) & provenance only (NOTRACE) \\
\midrule
100 & 6.7\% [2.2, 11.0] & 1.3\% [--3.8, 4.0] & 11.9\% [4.1, 15.5] \\
1,000 & 14.9\% [10.4, 17.6] & 1.6\% [--2.1, 4.8] & 14.3\% [10.2, 17.6] \\
5,000 & 17.4\% [13.9, 19.9] & 1.1\% [--2.5, 3.2] & 17.2\% [14.6, 20.0] \\
20,000 & 19.0\% [11.8, 28.3] & 2.3\% [--5.5, 8.3] & 16.3\% [8.6, 24.5] \\
\bottomrule
\end{tabular}
\end{table}

Three results, with their cross-campaign caveat stated first:
absolute times are not comparable between campaigns (baselines moved
with machine state across sessions); every comparison here is
\emph{within} one internally-controlled campaign, and the naive-vs-lazy
arc compares relative overheads, each measured against its own
contemporaneous baseline.

(1) \textbf{Lazy recording is statistically free}: 1--2\% at every size
with every CI spanning zero --- down from +143\%, a \textasciitilde{}60$\times$ reduction,
achieved with no loss of output (the trace-content tests pin
byte-identical details and passed unchanged, 205/205).
(2) \textbf{Provenance costs a steady 12--17\%} --- the honest correction of
the blocked campaigns' contradictory readings (--1.2\% and +19.9\% were
both drift). The cost is structural and now grounded, not
conjectured: adding \texttt{span} + \texttt{note} grew \texttt{sizeof(IRInstruction)}
from 116 to 160 bytes (+38\%; 20 bytes of span, 24 of MinGW's
\texttt{std::string}, measured), so every vector element, copy, and pass
traversal moves more data, and codegen additionally builds
provenance text per instruction. Reduction paths (a compressed span
representation; gating comment emission) are future work with a
target number attached, not hand-waving: the fields are
value-carried by design (D10), and that design's price is measured.
(3) At startup-amortized sizes the costs compose additively
(20k: 19.0 $\approx$ 2.3 + 16.3; 5k: 17.4 vs 18.3), consistent with
independent mechanisms; the 100-line row does \emph{not} compose (6.7 vs
13.2 summed) --- at 25 ms total runtime, process startup dominates and
the CIs are too wide for attribution, which is why the size ladder
exists. The 20k CIs are themselves wide ([11.8, 28.3] for FULL);
tightening them means more runs at the largest size, noted as a
cheap improvement rather than done, since no conclusion here turns
on the difference between 15\% and 20\%.

\textbf{Memory (E2).} Peak working set is unaffected throughout: $\leq$1.3\%
delta at every size in every campaign (interleaved, 20k: 85.8 MiB
baseline $\rightarrow$ 86.9 MiB full). Open-frame storage and per-instruction
fields are memory-negligible; the cost of provenance is time, not
space.

\subsection{Diagnostic quality (E3/E4)}

\textbf{E3 --- structural completeness (component presence, not quality).}
E3's epistemic role must be stated precisely: it is a structural
\emph{inventory}, not a comparison study. That our output contains the
sections we designed is near-tautological; the informative content
is the contrast \emph{shape} --- which explanation components a
conventional pipeline emits at all --- and the confirmation that no
component silently fails to render on any error class. Nine
single-error programs through both compilers on identical inputs;
mechanical detection rules only (the presence-rubric approach
descends from the error-message evaluation instrument of \cite{marceau2011measuring},
narrowed from judging responses to detecting components). Ours presents location,
cause, explanation prose, fix suggestions, trace, and caret on 9/9
(plus invalid/valid examples exactly on the two malformed-expression
cases, where the kind defines them); GCC 6.3.0 presents location,
cause, and caret on 9/9, and no note, fix-it, example, or trace on
any of the nine. Two scope statements keep this honest: the baseline
is the 2016 toolchain this compiler targets --- no claim is made
against current GCC/Clang without running them --- and presence is not
helpfulness; the enhanced-messages literature is contested on the
latter \cite{denny2014enhancing,pettit2017enhanced}, which is precisely why the paper measures the former.

\textbf{E4 --- cascade behavior under single-fault mutation.} 100 mutants,
each exactly one token-level edit from a valid seed (delete /
duplicate / undeclared-substitution at use sites only --- declaration
sites are excluded because corrupting a declaration creates k
\emph{genuine} downstream faults and would poison the one-fault ground
truth). Error diagnostics per rejected mutant:

\begin{table}[t]
\centering
\small
\begin{tabular}{llll}
\toprule
compiler & rejected & accepted & distribution \\
\midrule
ours & 100 & 0 & median 1, mean 1.32, max 5 (73$\times$1, 24$\times$2, 2$\times$3, 1$\times$5) \\
g++ 6.3 & 90 & 10 & median 1, mean 1.14, max 2 (77$\times$1, 13$\times$2) \\
\bottomrule
\end{tabular}
\end{table}

Findings, stated against ourselves first: (1) the shape-specific
cascade-freedom invariants do not generalize for free --- 3\% of mutants
still produce $\geq$3 diagnostics, so the tested-invariant approach is
necessary but not sufficient, and mutation testing is how the
residual frontier gets measured (the discovered shapes are enumerable
future fixes); (2) g++ is slightly tighter on this corpus; (3) the
accepted sets differ because the \emph{languages} differ (C permits \texttt{;;}
and \texttt{{{}; this language does not), so the distributions are not over
identical mutant sets --- an intersection-set statistic is listed as a
refinement in \S13; (4) all 100 mutants derive from a single 78-line
seed, so mutation diversity is bounded by one program's token
population --- a multi-seed corpus is the other \S13 refinement, and no
generalization beyond this corpus is claimed. Framing: cascading error reporting is a
quantified problem in the recovery literature \cite{diekmann2020dont}; our contribution
here is the \emph{property-test} treatment, not a recovery algorithm.

\textbf{E5 --- invasiveness.} From the introducing commit's diff, compiler
source only: trace recording $\approx$200 lines (recorder + three stage
integrations + surfacing), provenance $\approx$115 lines (IR fields,
generator stamping, two pass-preservation sites, assembly comments,
JSON export) --- $\approx$320 lines across 17 files, 5.6\% of compiler source,
supporting H1's "localized mechanisms" claim with a number. The
diagnostics subsystem overall was 28.2\% of compiler source at the
mechanism-introducing commit and 29.3\% at the time of writing ---
\S8.1 traces the growth --- the architecture's priorities, measured.

\subsection{The test suite as evidence}

205 tests (603 assertions) function as more than regression cover:
the trace-accuracy tests are the \emph{evidence mechanism} for H2 (they
assert runtime facts no static chain could contain, and their
failure at the curated$\rightarrow$recorded transition is documented evidence
the traces changed nature); the cascade tests pin the quality
invariants E4 then stress-tests; the byte-stability of golden IR
tests across the provenance change demonstrates the compatibility
cost D10 accepted; and the \texttt{--check} silence tests protect the
measurement channel every number in \S9.2 depends on.



